\textbf{Momento angular, producto vectorial y sistema de ecuaciones.} Primero se calculan las tres filas de $\boldsymbol H=I\boldsymbol\omega$. Las entradas $xz$ y $zx$ son negativas con el convenio del enunciado.\[\fitmath{\boldsymbol H=\left[\begin{matrix}1342.4\\740.0\\-1826.8\end{matrix}\right]\,\mathrm{kg\,m^2/s}}\]\[\fitmath{H_{1}=(9800)(0.13)+(-760)(-0.09)\approx1342.4\,\mathrm{kg\,m^2/s}}\]\[\fitmath{H_{2}=(14800)(0.05)\approx740\,\mathrm{kg\,m^2/s}}\]\[\fitmath{H_{3}=(-760)(0.13)+(19200)(-0.09)\approx-1826.8\,\mathrm{kg\,m^2/s}}\]\[\fitmath{c_{1}=(0.05)(-1826.8)-(-0.09)(740)\approx-24.74\,\mathrm{N\,m}}\]\[\fitmath{c_{2}=(-0.09)(1342.4)-(0.13)(-1826.8)\approx116.668\,\mathrm{N\,m}}\]\[\fitmath{c_{3}=(0.13)(740)-(0.05)(1342.4)\approx29.08\,\mathrm{N\,m}}\]El producto giroscópico se desarrolla como $(qH_z-rH_y,\ rH_x-pH_z,\ pH_y-qH_x)$; su valor es:\[\fitmath{\boldsymbol\omega\times\boldsymbol H=\left[\begin{matrix}-24.74\\116.668\\29.08\end{matrix}\right]\,\mathrm{N\,m}}\]Se resta de los momentos aplicados, obteniendo el segundo miembro $\boldsymbol b=\boldsymbol M-\boldsymbol\omega\times\boldsymbol H$.\[\fitmath{\boldsymbol b=\left[\begin{matrix}624.74\\-496.668\\390.92\end{matrix}\right]\,\mathrm{N\,m}}\]La ecuación de $y$ queda desacoplada: $\dot q=b_y/I_{yy}$. Para $x,z$, eliminar la otra aceleración en el sistema $I_{xx}\dot p-I_{xz}\dot r=b_x$, $-I_{xz}\dot p+I_{zz}\dot r=b_z$ da:\[\fitmath{\dot p=\frac{I_{zz}b_x+I_{xz}b_z}{I_{xx}I_{zz}-I_{xz}^2},\qquad\dot r=\frac{I_{xz}b_x+I_{xx}b_z}{I_{xx}I_{zz}-I_{xz}^2}}\]\[\fitmath{(\dot p,\dot q,\dot r)=\left[\begin{matrix}0.0655291072083522 & -0.0335586486486487 & 0.0229542771603306\end{matrix}\right]\,\mathrm{s^{-2}}}\]Para comprobar, multiplicar el vector de aceleraciones por $I$ y sumar el término giroscópico recupera $\boldsymbol M$. El trabajo giroscópico se anula por perpendicularidad; la tasa de energía es $\boldsymbol\omega\cdot\boldsymbol M=21.2\,\mathrm W$.
